% from ega2-1-fr.tex
\oldpage[II]{178}
qui est quasi-cohérente (\hyperref[I.9.2.2-fr]{I, 9.2.2, a)}). On a un
homomorphisme canonique de $\mathcal{O}_X$-Algèbres graduées
\begin{equation}
\label{II.8.8.1.2-fr}
  \tau:f^*(\mathcal{S})\longrightarrow
  \bigoplus_{n\geqslant 0}\mathcal{L}^{\otimes n}
\tag{8.8.1.2}
\end{equation}
qui, pour les degrés $\geqslant 1$, coïncide avec l'homomorphisme canonique
$\sigma:f^*(f_*(\mathcal{L}^{\otimes n}))\to\mathcal{L}^{\otimes n}$
(\hyperref[0.4.4.3-fr]{0, 4.4.3}), et pour le degré $0$ est l'identité.
L'hypothèse que $\mathcal{L}$ est $f$-ample entraîne alors
(\hyperref[II.4.6.3-fr]{4.6.3} et \hyperref[II.3.6.1-fr]{3.6.1}) que le
$Y$-morphisme correspondant
\begin{equation}
\label{II.8.8.1.3-fr}
  r=r_{\mathcal{L},\tau}:X\longrightarrow
  P=\operatorname{Proj}(\mathcal{S})
\tag{8.8.1.3}
\end{equation}
est partout défini et est une \emph{immersion ouverte dominante}, et que l'on
a
\begin{equation}
\label{II.8.8.1.4-fr}
  r^*(\mathcal{O}_P(n))=\mathcal{L}^{\otimes n}
  \qquad\text{pour tout }n\in\mathbf{Z}.
\tag{8.8.1.4}
\end{equation}
\end{env}

\begin{proposition}[8.8.2]
\label{II.8.8.2-fr}
Soit $C=\operatorname{Spec}(\mathcal{S})$ \emph{le cône affine défini par
$\mathcal{S}$}~; si $\mathcal{L}$ est $f$-ample, il existe un $Y$-morphisme
canonique
\begin{equation}
\label{II.8.8.2.1-fr}
  g:V=\mathbf{V}(\mathcal{L})\longrightarrow C
\tag{8.8.2.1}
\end{equation}
tel que le diagramme
\begin{equation}
\label{II.8.8.2.2-fr}
  \xymatrix{
    X\ar[r]^{j}\ar[d]_{f} &
    \mathbf{V}(\mathcal{L})\ar[r]^{\pi}\ar[d]^{g} & X\ar[d]^{f}\\
    Y\ar[r]_{\varepsilon} & C\ar[r]_{\psi} & Y
  }
\tag{8.8.2.2}
\end{equation}
soit commutatif, $\psi$ et $\pi$ étant les morphismes structuraux, $j$ et
$\varepsilon$ les immersions canoniques appliquant respectivement $X$ et $Y$
sur la section nulle de $\mathbf{V}(\mathcal{L})$ et le préschéma des sommets
de $C$. En outre, la restriction de $g$ à
$\mathbf{V}(\mathcal{L})-j(X)$ est une immersion ouverte
\begin{equation}
\label{II.8.8.2.3-fr}
  \mathbf{V}(\mathcal{L})-j(X)\longrightarrow
  E=C-\varepsilon(Y)
\tag{8.8.2.3}
\end{equation}
dans le cône épointé affine $E$ correspondant à $\mathcal{S}$.
\end{proposition}

\begin{proof}
Avec les notations de (\hyperref[II.8.8.1-fr]{8.8.1}), soient
$\mathcal{S}_P^{\geqslant}=\bigoplus_{n\geqslant 0}\mathcal{O}_P(n)$, et
$C_P=\operatorname{Spec}(\mathcal{S}_P^{\geqslant})$. On sait
(\hyperref[II.8.6.2-fr]{8.6.2}) qu'on a un morphisme canonique
$h=\operatorname{Spec}(\alpha):C_P\to C$ tel que le diagramme
\begin{equation}
\label{II.8.8.2.4-fr}
  \xymatrix{
    C_P\ar[r]\ar[d]_{h} & P\ar[d]^{p}\\
    C\ar[r]_{\psi} & Y
  }
\tag{8.8.2.4}
\end{equation}
soit commutatif~; en outre, si $\varepsilon_P:P\to C_P$ est l'immersion
canonique, le diagramme
\begin{equation}
\label{II.8.8.2.5-fr}
  \xymatrix{
    P\ar[r]^{\varepsilon_P}\ar[d]_{p} & C_P\ar[d]^{h}\\
    Y\ar[r]_{\varepsilon} & C
  }
\tag{8.8.2.5}
\end{equation}
est commutatif (\hyperref[II.8.7.1.1-fr]{8.7.1.1}), et enfin, la restriction
de $h$ au cône épointé affine $E_P$ est un \emph{isomorphisme}
$E_P\xrightarrow{\sim}E$ (\hyperref[II.8.6.2-fr]{8.6.2}). D'autre part, il
résulte de (\hyperref[II.8.8.1.4-fr]{8.8.1.4}) que l'on a
\[
  r^*(\mathcal{S}_P^{\geqslant})=
  \mathbf{S}_{\mathcal{O}_X}(\mathcal{L})
\]

